\documentclass[a4paper]{article}

\usepackage[fontset=fandol]{ctex}
\usepackage{amsmath,amssymb}
\usepackage{graphicx}
\usepackage{xcolor}
\usepackage[margin=2.2cm]{geometry}
\usepackage[most]{tcolorbox}
\usepackage{listings}
\usepackage{hyperref}
\usepackage{booktabs}
\usepackage{float}

\newtcolorbox{knowledgebox}[1]{
  enhanced, colback=blue!5!white, colframe=blue!75!black, colbacktitle=blue!75!black,
  coltitle=white, fonttitle=\bfseries, title=#1, sharp corners
}

\newtcolorbox{importantbox}[1]{
  enhanced, colback=yellow!10!white, colframe=yellow!80!black, colbacktitle=yellow!80!black,
  coltitle=black, fonttitle=\bfseries, title=#1, sharp corners
}

\newtcolorbox{warningbox}[1]{
  enhanced, colback=red!5!white, colframe=red!75!black, colbacktitle=red!75!black,
  coltitle=white, fonttitle=\bfseries, title=#1, sharp corners
}

\lstset{
  basicstyle=\ttfamily\small,
  breaklines=true,
  frame=single,
  numbers=left,
  numberstyle=\tiny\color{gray},
  captionpos=b
}

\newcommand{\notetitle}{UCB CS294/194-196 Agentic AI (Fall 2025) 第05讲：Agent Evaluation \& Project Overview}
\newcommand{\noteauthors}{Codex}
\newcommand{\notedate}{\today}
\newcommand{\videochannel}{Berkeley RDI / Agentic AI MOOC}
\newcommand{\videopublishdate}{2025-10-06}
\newcommand{\videoduration}{1:34:22}
\newcommand{\videourl}{https://www.youtube.com/watch?v=VfOA2a0dj4w}
\newcommand{\videocoverpath}{../../raw/05_VfOA2a0dj4w/05_VfOA2a0dj4w.jpg}

\begin{document}

\begin{titlepage}
\centering
{\Large 课程笔记\par}
\vspace{1.0cm}
{\huge\bfseries \notetitle\par}
\vspace{0.6cm}
{\large \noteauthors\par}
\vspace{0.25cm}
{\large \notedate\par}
\vspace{0.9cm}
\includegraphics[width=0.82\textwidth,height=0.45\textheight,keepaspectratio]{\videocoverpath}\par
\vfill
\begin{tcolorbox}[width=0.9\textwidth, colback=black!2!white, colframe=black!60, sharp corners]
\textbf{视频作者/频道}：\videochannel\par
\textbf{发布日期}：\videopublishdate\par
\textbf{视频时长}：\videoduration\par
\textbf{视频链接}：
\url{\videourl}
\end{tcolorbox}
\end{titlepage}

\tableofcontents
\newpage

\section{为什么 agent 评估是核心问题}
\subsection{评估不是附属品}
这一讲的第一句定义非常直接：evaluation 是对模型和 agent 的 \emph{systematic, repeatable measurement}。对于 agent 来说，评估的重要性比普通 LLM 更高，因为 agent 不仅要回答问题，还要在动态环境里规划、调用工具、维持记忆，并完成多步任务。没有可信评估，训练进步、部署安全和 benchmark 竞争都无从谈起。

\begin{table}[H]
\centering
\small
\begin{tabular}{@{}p{0.22\textwidth}p{0.33\textwidth}p{0.33\textwidth}@{}}
\toprule
维度 & LLM & Agent \\
\midrule
状态 & 静态文本输入输出 & 动态环境 + 工具 + 历史交互 \\
能力 & 生成、补全、总结 & 规划、执行、迭代、修复 \\
评估难点 & 文本质量与事实性 & 任务成功率、路径可靠性、工具使用、环境状态变化 \\
风险 & 幻觉、风格偏差 & 错误行动、状态污染、prompt/benchmark gaming \\
\bottomrule
\end{tabular}
\caption{Agent eval 相比普通 LLM eval 的复杂度提升。}
\end{table}

\subsection{什么时候必须做评估}
讲者给出的答案是：几乎在训练和上线的每个阶段都需要评估。训练前要评估数据质量和任务定义；训练中要看 loss、reward、pass rate 和中间行为；训练后要用 benchmark 和 harness 比较模型；上线前要看安全性、成本和用户体验；上线后还要持续监控漂移。对 agent 而言，评估不是最后一步，而是整个流水线里的闭环控制信号。

\begin{importantbox}{评价的三个用途}
1. \textbf{比较模型}：公平比较不同模型和不同策略。\\
2. \textbf{指导决策}：决定是否可以部署、是否需要回滚。\\
3. \textbf{推动进步}：让 benchmark 本身成为研究对象，而不是静态分数板。
\end{importantbox}

\subsection{本章小结}
这一节把评估的地位讲得很清楚：agentic 系统的每一次能力提升，最后都要通过可复现的评估来被确认。没有 evaluation，后续的训练、部署和 project grading 都失去参照物。

\section{评估的类型}
\subsection{close-ended, open-ended, verifiable, non-verifiable}
讲座把 agent 评估的第一层分类拆成几个很实用的维度。close-ended 任务的正确答案数量有限，因此容易自动化评估；open-ended 任务则更接近真实应用，但又必须进一步区分 \emph{verifiable} 与 \emph{non-verifiable}。前者像数学证明和代码生成，后者像故事创作、风格迁移和开放式写作。

\begin{table}[H]
\centering
\small
\begin{tabular}{@{}p{0.16\textwidth}p{0.22\textwidth}p{0.24\textwidth}p{0.26\textwidth}@{}}
\toprule
类别 & 典型任务 & 评价方式 & 代表问题 \\
\midrule
Close-ended & 多选、固定答案题 & 自动判分 & 答案空间是否有限？ \\
Open-ended / verifiable & 数学、代码 & oracle、测试、checker & 是否存在明确正确性标准？ \\
Open-ended / non-verifiable & 写作、风格、摘要 & 人类评估、LLM judge & 结果是否高度主观？ \\
Static & 固定测试集 & 可复现对比 & 会不会很快被过拟合或污染？ \\
Dynamic & 持续更新测试集 & 持续追踪 & 能否保持相关性与抗污染？ \\
\bottomrule
\end{tabular}
\caption{讲座中的评估分类。}
\end{table}

\subsection{人类评估与 LLM-as-a-Judge}
人类评估通常被视为 gold standard，但它慢、贵、不可重复，而且不同标注者之间可能存在 disagreement，甚至同一标注者在不同时刻也会给出不同判断。LLM-as-a-judge 则提高了吞吐，但会遇到随机性、偏置、分数难解释、对 prompt 很敏感等问题。讲者的建议不是简单放弃，而是把它们当作可用但不完美的信号：要和 human eval 交叉验证，要做多次重复打分，要尽量用更细的 rubric 和连续分数。

\subsection{静态评估与动态评估}
静态 benchmark 的优点是稳定、可复现，适合学术比较；动态 benchmark 则更接近真实世界数据分布，也更不容易被污染或过拟合。举例来说，ImageNet、GLUE、MMLU 属于典型静态 benchmark，而 DynaBench、LiveCodeBench 则更强调动态更新。对于 agent 来说，动态环境几乎是必需的，因为工具、网页、API 和用户需求本来就在变化。

\subsection{本章小结}
这一部分给了评估的基础 taxonomy。最重要的不是名词，而是判断：当任务有明确 oracle 时，应该尽量用 verifiable eval；当任务难以自动判分时，必须小心设计 human/LLM rubric，并警惕静态 benchmark 很快被污染或饱和。

\section{什么是好的 eval system}
\subsection{Outcome validity}
讲座用一个关键词概括“好评估系统”应该追求的东西：\emph{outcome validity}。它的意思不是看模型是不是在某个 benchmark 上拿了高分，而是看这个分数是否真的对应了我们关心的真实任务结果。比如，在文本、代码、环境状态变化和多步推理上，判分信号都应尽量对应真正的 task success。

\begin{knowledgebox}{Outcome validity 的直觉}
一个好 eval 不应该只测“模型会不会讨巧”，而应该测“模型完成真实任务的能力是否真的更强了”。如果评估可以被 shortcut、gaming 或数据污染轻易击穿，那么它就失去了定义 intelligence 的资格。
\end{knowledgebox}

\subsection{评估最常见的失败方式}
讲者专门列出了几个常见坑：数据噪声和偏置会扭曲结论；评估如果不切实际，就无法反映 practitioner 的真实需求；如果 benchmark 可以被 shortcut，就会鼓励投机；如果任务不够难，就会很快饱和；如果测试集过于静态，公开之后就会污染得很快。简言之，评估既要可复现，又要难作弊，还要足够贴近真实世界。

\begin{warningbox}{一个很实用的判断标准}
如果你发现模型在 benchmark 上不断变强，但真实任务效果没有同步提升，那就该怀疑 evaluation 本身是否已经失真。
\end{warningbox}

\subsection{本章小结}
好的 eval system 不是“能算分”这么简单，而是要有 outcome validity。它要能解释真实任务成功、支持鲁棒比较，并且尽量抵抗污染、捷径和过拟合。

\section{案例研究：好 benchmark 是怎么做出来的}
\subsection{CyberGym, \texorpdfstring{$\tau$}{tau}-bench, \texorpdfstring{$\tau^2$}{tau2}-bench}
这一讲用多个案例说明如何把一个 benchmark 从“看起来像任务”做成真正可用的评估系统。CyberGym 的数据来自 ARVO 和 OSS-Fuzz 等真实世界漏洞来源，通过重建 pre/post patch commits、可执行文件和 ground-truth PoC，再用 runtime sanitizers 做自动判定。它的关键点是：不靠主观评分，而是通过是否真正触发漏洞来判定成功。

\(\tau\)-bench 则把重点放在 tool-agent-user interaction 上。任务通常是多轮对话式的，agent 要在遵守 policy 的前提下，调用 API 工具完成用户目标，最后通过数据库状态与目标状态是否一致来判断成败。它天然支持 pass@1 与 pass@k 这类可靠性指标。

\(\tau^2\)-bench 在 \(\tau\)-bench 的基础上进一步系统化，加入 LLM-drafted PRD、mock DB、unit tests，以及环境、通信、自然语言和 action 等多种 assert 分类。它想表达的核心是：benchmark 不只是题目集合，而是一套能自动验证的任务生成管线。

\subsection{GDPval, CRMArena, LegalAgentBench}
GDPval 的特点是把任务交给有丰富经验的专业人士来写，再由同领域的专家盲审 AI 与 human deliverables。它特别强调低 contamination 风险和真实工作流，适合衡量“像人一样完成工作”的能力。CRMArena 则更偏向 CRM 业务流程里的多步代理任务。LegalAgentBench 则把法律场景中的多跳检索、逻辑推理和法律写作整合到一个多工具环境里，工具数量和语料数量都相当大。

\begin{table}[H]
\centering
\small
\begin{tabular}{@{}p{0.18\textwidth}p{0.25\textwidth}p{0.2\textwidth}p{0.28\textwidth}@{}}
\toprule
Benchmark & 任务域 & 判分方式 & 关键设计点 \\
\midrule
CyberGym & cybersecurity & runtime sanitizer / crash trace & 真实漏洞、pre/post patch、PoC \\
\(\tau\)-bench & tool-agent-user & final DB state vs goal state & policy following、pass@k \\
\(\tau^2\)-bench & complex workflows & categorical assertions & PRD、mock DB、compositional tasks \\
GDPval & professional work & expert blind grading & 低污染、真实 deliverable \\
LegalAgentBench & legal workflows & multi-tool success & 17 corpora、37 tools、多跳推理 \\
\bottomrule
\end{tabular}
\caption{讲座中用来说明“好 benchmark 应该长什么样”的案例。}
\end{table}

\subsection{本章小结}
这些案例共同说明了一件事：好 benchmark 不是简单堆题，而是把 task、environment、metric、data pipeline 和 validation 一起设计好。越接近真实任务，越要重视 low contamination、difficulty stratification 和可自动验证的终点条件。

\section{课程项目：green agent 应该怎么做}
\subsection{green agent 与 white agent}
课程项目采用了一个很形象的分工：\emph{green agent} 是评估方或主持方，\emph{white agent} 是被测试的系统。green agent 的职责是 orchestrate 整个评估流程，读取任务、调用工具、收集结果、做判分，并把最终评估结果写回平台。它本质上就是一个评估技术员（technician），不是一个单纯跑分脚本。

\subsection{两类项目}
讲义把 project 分成两类。第一类是 \emph{integrating existing benchmarks}：把已经发表、验证过的 benchmark 接到 AgentBeats / 课程平台里，再对其质量做分析、修正和扩展。第二类是 \emph{building new benchmarks}：围绕一个新的任务域，自己设计环境、数据管线、评估指标和 green agent。

\begin{table}[H]
\centering
\small
\begin{tabular}{@{}p{0.2\textwidth}p{0.34\textwidth}p{0.34\textwidth}@{}}
\toprule
类型 & 主要工作 & 课程里特别看重什么 \\
\midrule
Integrate existing benchmark & 适配已有 benchmark、分析质量问题、修正与扩展 & Faithfulness、质量分析、可复用性 \\
Build new benchmark & 设计新任务、环境、metric、test cases、grader & Novelty、scope、validation、realism \\
\bottomrule
\end{tabular}
\caption{课程项目的两种主路线。}
\end{table}

\subsection{Step-by-step checklist}
讲座给出的 checklist 非常实用：先选任务，再设计环境，再设计 metrics，最后设计 test cases。以订票 agent 为例，环境要定义浏览器或 APP、动作空间、环境反馈；metric 要定义成功率、价格、是否满足用户要求；test cases 要覆盖成功、失败、错误预订、找不到网站等多种情形。对 existing benchmark 来说，还要补上 manual validation、evaluator check、bias/limitation notes，并在必要时做 correction 和 expansion。

\subsection{Project grading rubric}
对新 benchmark，rubric 重点看 goal \& novelty、scope \& scale、evaluator quality、validation、reliability、quality assurance、realism 和 impact。对已有 benchmark，rubric 重点看 analysis、faithfulness、quality assurance、evaluator quality、validation、reliability 和 impact。换句话说，老师并不只看“你能不能把 benchmark 跑起来”，还看你是否理解它为什么成立、哪里有缺陷、你修复后是否更可靠。

\subsection{本章小结}
项目部分把整门课的评估思想落到可执行任务上：做 green agent 不是只做一个 wrapper，而是要完整设计任务、环境、metric、测试与判分逻辑。对于新 benchmark，重点是新能力空间；对于已有 benchmark，重点是质量分析和修复能力。

\section{总结与延伸}
这一讲把 agent eval 从“刷分”提升成了一门完整的方法论。最核心的结论是：评估必须和真实任务的 outcome 对齐，必须尽量可复现，并且必须能抵抗污染、捷径和饱和。对 agent 来说，评估不只是最后的验收，而是训练、部署和课程项目本身的一部分。

如果把这一讲和前后的内容连起来看，就会得到一个很清晰的 pipeline：先用可验证训练数据和奖励训练 agent，再用合适的 benchmark 与 harness 评估 agent，最后把评估机制本身也做成可以持续改进的系统。对于后续做课程项目的同学来说，真正的高分不在于“把模型跑出来”，而在于“把评估做对”。

\end{document}
